% --- Back cover ---
% Large PERSONIX wordmark, no logo image. Cream paper inherited globally.
% Absolute (tikz overlay) positioning so the footer sticks to the bottom.
\clearpage
\thispagestyle{empty}
\begin{tikzpicture}[remember picture, overlay]
  % --- Wordmark ---
  \node[anchor=north, font=\sffamily\bfseries\fontsize{48}{52}\selectfont]
    at ([yshift=-26mm]current page.north) {PERSONIX};
  \node[anchor=north, font=\rmfamily\itshape\large]
    at ([yshift=-44mm]current page.north) {समझौता-रहित परिवर्तन};
  \node[anchor=north, font=\rmfamily\small,
        text width=\paperwidth-50mm, align=center]
    at ([yshift=-54mm]current page.north)
    {सेंसरशिप-रोधी और भ्रष्टाचार-रोधी विकेंद्रीकृत प्रतिष्ठा नेटवर्क};
  % --- Body ---
  \node[anchor=north, text width=\paperwidth-50mm, align=justify, font=\small]
    at ([yshift=-72mm]current page.north)
    {%
      राज्य तब काम करता था जब संचार धीमा था, निर्णय स्थानीय थे, और सत्ता उसी
      शहर में रहती थी जिस पर वह शासन करती थी। आज के नेटवर्क तीनों धारणाओं को
      उलट देते हैं --- और उन पर निर्मित संस्थाएँ अब उस संसार से मेल नहीं खातीं
      जिस पर वे शासन करती हैं। यह पुस्तक एक ऐसे उपकरण का वर्णन करती है जो मेल
      खाता है: एक विकेंद्रीकृत प्रतिष्ठा नेटवर्क जिसमें पहचान आपकी अपनी रहती है,
      सत्य सार्वजनिक रूप से लेखा-परीक्षण-योग्य है, और रिश्वतखोरी प्रतिष्ठा की
      आत्महत्या है।

      \smallskip
      न कोई घोषणापत्र। न कोई भविष्यवाणी। इस बात का एक कार्यशील खाका कि राज्य
      का उत्तराधिकारी कैसे बनाया जा सकता है --- स्वेच्छा से, एक-एक घोषणा करके।
    };
  % --- Footer (anchored to the bottom of the page) ---
  \node[anchor=south, font=\sffamily\footnotesize,
        text width=\paperwidth-50mm, align=center]
    at ([yshift=12mm]current page.south)
    {%
      Pavel Kudrna\quad\textbullet\quad प्राग, 2026\par
      \href{https://personix.org}{personix.org}\par
      Creative Commons Attribution-ShareAlike 4.0 International
      (CC BY-SA 4.0) के अंतर्गत लाइसेंस-प्राप्त।
    };
\end{tikzpicture}%
\mbox{}%
\clearpage
